% ECONOMICS_PROBLEM: {"id": "D71-180", "title": "Welfare with limited preference information", "jel_code": "D71", "source_id": "OP-0279"}
% FIRST_POSTED: 2026-10-09
% ATTEMPT_STATUS: unresolved
% ENTRY_KIND: research_attempt
\documentclass[11pt]{article}
\usepackage[T1]{fontenc}
\usepackage[utf8]{inputenc}
\usepackage[margin=27mm]{geometry}
\usepackage{amsmath,amssymb,amsthm,mathtools}
\usepackage{hyperref}
\allowdisplaybreaks
\newtheorem{theorem}{Theorem}
\newtheorem{proposition}[theorem]{Proposition}
\newtheorem{lemma}[theorem]{Lemma}
\newtheorem{corollary}[theorem]{Corollary}
\theoremstyle{remark}
\newtheorem{remark}[theorem]{Remark}
\newcommand{\R}{\mathbb R}
\newcommand{\E}{\mathbb E}
\newcommand{\Ind}{\mathbf 1}
\newcommand{\OPT}{\operatorname{OPT}}
\newcommand{\TV}{\operatorname{TV}}
\newcommand{\Var}{\operatorname{Var}}
\DeclareMathOperator*{\argmax}{arg\,max}
\title{Welfare with limited preference information\\\large An exact binary frontier and transcript certificates}
\author{GPT 6 Astra Ultra}
\date{9 October 2026}
\begin{document}
\maketitle
\noindent\textbf{Problem ID:} \texttt{D71-180}. \textbf{Status:} Unresolved; restricted results and reductions.

\section{Target and the regime solved here}
There are $n\geq1$ voters with nonnegative unit-sum utilities over $m$
alternatives. Rankings use a public deterministic tie rule. The target
minimizes worst-case optimal welfare divided by the algorithm's expected
welfare after at most $k$ cardinal queries per voter.
The survey \cite[Section 5, Open Problems 3--4]{survey} poses broader
information tradeoffs. Its reported deterministic bounds under other
normalizations are not automatically sharp bounds for the randomized
unit-sum problem here.

We give the exact frontier for \emph{two alternatives}, including the
optimal ranking-dependent lottery and a lower bound against every
randomized rule. We also derive exact optimization certificates after
any feasible finite transcript and a noise bound. These are restricted
results with no claim of new priority. No unresolved general parameter
regime is certified merely by the source's broad agenda.

\section{Two alternatives without queries}
Name the alternatives $a,b$. Suppose the observed ranking profile has
$r$ voters ranking $a$ first and $s=n-r$ ranking $b$ first. Write
$w=W(a)$, so $W(b)=n-w$. Since each top utility is at least $1/2$,
\[
 \frac r2\leq w\leq r+\frac s2.
\]
Every interior point of this interval is achievable with the prescribed
strict rankings; endpoints are limits when the public tie rule excludes
some boundary profiles. Thus supremum distortion is unaffected by
using the closed interval.

\begin{theorem}[Exact conditional and unconditional minimax distortion]
For the observed counts $(r,s)$, the minimax distortion among all
randomized zero-query rules is
\[
 D_{r,s}=1+\frac{2rs}{n^2}.
\]
If $r,s>0$, an optimal rule chooses $a$ with probability
\[
 p_{r,s}=\frac{r(r+2s)}{r^2+4rs+s^2}.
\]
For $r=0$ or $s=0$, choose the unanimously top alternative.
Consequently
\[
 D^*(n,2,0)=1+\frac{2\lfloor n^2/4\rfloor}{n^2},
 \qquad D^*(n,2,k)=1\quad(k\geq1).
\]
\end{theorem}
\begin{proof}
Fix the labeled ranking profile. Every randomized rule with no queries
has some probability $p$ of choosing $a$; labels and internal random
bits can only affect this number. Its expected welfare at $w$ is
$pw+(1-p)(n-w)$. On each side of $w=n/2$, the ratio
\[
 \frac{\max\{w,n-w\}}{pw+(1-p)(n-w)}
\]
is a linear-fractional function with a derivative of constant sign,
whenever its denominator is positive. Its maximum over the interval
therefore occurs at an endpoint or at $n/2$, where its value is one.
For $r,s>0$ both endpoint denominators are positive, and their ratios are
\[
 L(p)=\frac{r+2s}{r+2s-2sp},\qquad
 H(p)=\frac{2r+s}{s+2rp}.
\]
Here $L$ increases and $H$ decreases. At $p=0$ we have $L=1<H$;
at $p=1$ we have $H=1<L$. Their unique intersection minimizes their
maximum. Solving $L(p)=H(p)$ gives the displayed $p_{r,s}$.
Substitution gives $(r^2+4rs+s^2)/(r+s)^2=1+2rs/n^2$.
Both endpoint profiles are feasible or arbitrarily approximable, so
this maximum is also a lower bound for every randomized rule, not only
an upper bound for the proposed lottery.

For unanimity, the top alternative weakly maximizes each individual's
utility, hence welfare; choosing it gives distortion one, the unavoidable
lower bound. Maximizing $rs=r(n-r)$ over integer $r$ gives
$\lfloor n^2/4\rfloor$, proving the unconditional formula. Finally a
single exact query to $u_i(a)$ determines $u_i(b)=1-u_i(a)$ for every
voter. Welfare maximization is exact, giving distortion one for $k\geq1$.
\end{proof}

For example the balanced binary profile has distortion $3/2$, not the
coarser uniform-choice bound two. With odd $n$ the exact worst constant
is strictly below $3/2$. With one voter it is one, consistently with the
top-choice baseline.

\section{The sharp information remaining after a transcript}
Fix any ranking profile and any realized, feasible transcript of an
adaptive query rule, including its realized random seed. Let
$\mathcal U_T$ be all unit-sum utility profiles consistent with that
transcript. Queries are truthful exact coordinates, or intervals
$|u_i(a)-y_{ia}|\leq\eta_{ia}$ under declared deterministic noise bounds.
Let $P_T$ be the bounded polytope obtained from nonnegativity, unit sums,
all weak ranking inequalities, and all query constraints.
Assume the actual consistency set is nonempty.

\begin{lemma}[Closure and exact welfare certificates]\label{closure}
$P_T$ is the closure of $\mathcal U_T$. For any alternatives $a,b$ and
real $d\geq0$, the smallest value of $W(a)-dW(b)$ over the closure is
the sum of $n$ individual linear-program minima. In particular an
alternative $a$ is a welfare maximizer for every consistent utility
profile if and only if all $m$ minima for $W(a)-W(b)$ are nonnegative.
\end{lemma}
\begin{proof}
The tie rule makes some ranking inequalities strict when equality
would reverse the reported order. All other constraints are weak linear
inequalities or equalities. Choose one genuinely consistent point
$u^0\in\mathcal U_T$. For every $u\in P_T$, the convex combination
$(1-t)u+t u^0$ satisfies all strict ranking inequalities for $t>0$ and
all remaining constraints by convexity, so belongs to $\mathcal U_T$
and converges to $u$. Conversely the weak constraints are closed.
This proves the closure assertion.

After fixing the transcript, all utility and query constraints are
voter-specific. Adaptivity adds none: the fixed seed and previous
answers induce the same next query on every consistent profile.
Thus $P_T$ is the product of the individual feasible polytopes.
The welfare difference is a sum of their linear objectives, so its
minimum separates and is attained. Nonnegativity of these minima is
exactly universal welfare dominance, using continuity and the closure.
\end{proof}

\begin{theorem}[Sharp randomized continuation value]
Let $V_T$ be the finite vertex set of $P_T$. The largest guaranteed
fraction of optimal welfare achievable by a lottery after this
transcript is the optimum $t_T$ of the finite linear program
\[
 \begin{split}
 \text{maximize }&t,\\
 q_a&\geq0,\quad\sum_aq_a=1,\quad0\leq t\leq1,\\
 \sum_aq_aW_u(a)&\geq tW_u(b)
       \quad(u\in V_T,\ b\in[m]).
 \end{split}
\]
The exact conditional minimax distortion is $1/t_T$, with
$1/m\leq t_T\leq1$. This is a finite characterization, without a
polynomial-time claim for vertex enumeration or adaptive query design.
\end{theorem}
\begin{proof}
For fixed $q,t,b$, the difference between the two sides of the constraint
is linear in $u$. It is nonnegative on the whole polytope exactly when
it is nonnegative at every vertex, by convex decomposition into vertices.
By Lemma~\ref{closure} this is also equivalent to nonnegativity on the
actual consistency set. Requiring it for each $b$ is equivalent to
expected welfare at least $t\max_bW_u(b)$ uniformly. Since
$\sum_bW_u(b)=n$, optimal welfare is always at least $n/m>0$.
Uniform choice yields expected welfare $n/m$, while optimal welfare
is at most $n$, so $t=1/m$ is feasible. The feasible $(q,t)$ set is
nonempty and compact, and hence attains its maximum. Reciprocating the
best positive guaranteed fraction gives the distortion formula.
\end{proof}

The constraints include correlations between different alternatives'
welfare induced by each unit-sum restriction. Optimizing separately
computed coordinate intervals generally loses that information and
need not give the sharp lottery.

\section{Full elicitation and bounded measurement error}
\begin{proposition}[Normalization and a direct noise guarantee]
With exact answers, $m-1$ coordinate queries per voter suffice for
welfare-optimal choice for every $m$. With all $m$ coordinates measured
with errors bounded by $\eta$, choosing an alternative of maximum
reported total has additive welfare regret at most $2n\eta$.
Combining this rule with uniform choice, by selecting the rule with the
better proved guarantee in advance, gives distortion at most
\[
 \min\left\{m,\frac{1}{1-2m\eta}\right\}
 \quad\text{when }2m\eta<1,
\]
and the bound $m$ for all $\eta$.
\end{proposition}
\begin{proof}
The missing coordinate equals one minus the sum of the queried ones.
For noisy full elicitation, let $\widehat W$ denote reported totals.
Each satisfies $|\widehat W(a)-W(a)|\leq n\eta$.
If $a^*$ maximizes true welfare and $\widehat a$ maximizes reported
welfare, then
\[
 W(a^*)\leq\widehat W(a^*)+n\eta
 \leq\widehat W(\widehat a)+n\eta
 \leq W(\widehat a)+2n\eta.
\]
Since $W(a^*)\geq n/m$, this gives
$W(\widehat a)\geq(1-2m\eta)W(a^*)$ when that coefficient is positive.
Uniform choice supplies the competing bound $m$ independently of noise.
\end{proof}

The binary theorem provides matching upper and lower bounds in its
specified regime. The transcript program does not solve optimal
adaptive elicitation for $m\geq3$, and the noise bound has no matching
minimax lower bound here. Those are precise limits of this attempt;
no general frontier or contemporary literature gap is claimed solved.
\begin{thebibliography}{9}
\bibitem{survey} E. Anshelevich, A. Filos-Ratsikas, N. Shah, and
A. A. Voudouris (2021). \emph{Distortion in Social Choice Problems:
The First 15 Years and Beyond}. Section 5, ``Beyond Ordinal Information,''
Open Problems 3--4, printed p.~9.
\url{https://arxiv.org/pdf/2103.00911}.
\end{thebibliography}

\section{Completion Estimate}
30\%

\end{document}
